\documentclass[10pt]{article}
\usepackage[T1]{fontenc}
\usepackage[utf8]{inputenc}
\usepackage[french]{babel}
\usepackage{amsmath,amssymb,mathrsfs,mathpazo}
\usepackage{geometry}
\geometry{a4paper,margin=2cm}
\usepackage{graphicx,array,tabularx,tabularray,colortbl,multirow,multicol,booktabs,enumitem,pifont,ulem,fancybox,diagbox,cancel}
\usepackage{pstricks,pstricks-add,pst-plot,pst-tree,pst-node,pst-text,pst-func,pst-math,pst-bspline,pst-3dplot,pst-eucl,pst-solides3d}
\let\pscalculate\undefined
\usepackage{pst-bezier}
\usepackage{tikz,pgfplots,tkz-tab}
\usetikzlibrary{arrows,arrows.meta,patterns,calc,positioning,trees,decorations.markings,shapes.geometric,shapes.misc,shapes.arrows,intersections,decorations.pathmorphing,decorations.pathreplacing,backgrounds,shadows,babel}
\pgfplotsset{compat=1.18}
\usepgfplotslibrary{fillbetween}
\usepackage[np]{numprint}
\usepackage{listings,makecell,pas-tableur,twoopt}
\let\barre\undefined
\usepackage[pstricks]{tablvar}
\definecolor{darkgreen}{RGB}{0,100,0}
\providecommand{\red}{\color{red}}
\providecommand{\blue}{\color{blue}}
\providecommand{\green}{\color{green}}
\providecommand{\black}{\color{black}}
\providecommand{\white}{\color{white}}
\providecommand{\gray}{\color{gray}}
\providecommand{\cyan}{\color{cyan}}
\providecommand{\magenta}{\color{magenta}}
\providecommand{\yellow}{\color{yellow}}
\providecommand{\violet}{\color{violet}}
\providecommand{\orange}{\color{orange}}
\providecommand{\euro}{\texteuro}
\providecommand{\EUR}{\texteuro}
\providecommand{\ieme}{\textsuperscript{e}}
\providecommand{\degres}{\ensuremath{^\circ}}
\providecommand{\degree}{\ensuremath{^\circ}}
\providecommand{\vv}[1]{\overrightarrow{#1}}
\providecommand{\Coord}[1]{#1}
\providecommand{\decosix}{}
\providecommand{\starredbullet}{\ensuremath{\bullet}}
\pagestyle{empty}
\setlength{\parindent}{0pt}
\usepackage{adjustbox}\DeclareUnicodeCharacter{2236}{:}\DeclareUnicodeCharacter{2190}{\ensuremath{\leftarrow}}
\providecommand{\R}{\mathbb{R}}
\providecommand{\N}{\mathbb{N}}
\providecommand{\Z}{\mathbb{Z}}
\providecommand{\Q}{\mathbb{Q}}
\providecommand{\D}{\mathbb{D}}
\providecommand{\C}{\mathbb{C}}
\providecommand{\e}{{\mathrm{e}}}
\providecommand{\vect}[1]{\overrightarrow{#1}}
\providecommand{\vectt}[1]{\overrightarrow{\mathrm{#1}}}
\providecommand{\vvt}[1]{\overrightarrow{\mathrm{#1}}}
\providecommand{\np}[2][]{{#2~#1}}
\providecommand{\vv}[1]{{\overrightarrow{#1}}}
\renewcommand{\d}{{\mathrm{d}}}
\providecommand{\degres}{{^\circ}}
\providecommand{\degre}{{^\circ}}
\providecommand{\blue}{\color{blue}}
\providecommand{\red}{\color{red}}
\providecommand{\enskip}{\quad}
\providecommand{\llbracket}{\lbrack\!\lbrack}
\providecommand{\rrbracket}{\rbrack\!\rbrack}
\providecommand{\euro}{{\text{€}}}
\providecommand{\no}{{\text{n°}}}
\providecommand{\interval}[2]{{[#1;#2]}}
\providecommand{\intervFO}[2]{{[#1;#2[}}
\providecommand{\dotfill}{\ldots}

\providecommand{\ds}{\displaystyle}
\providecommand{\textding}[1]{#1}
\providecommand{\Oij}{\(\mathrm{O};\vec{\imath},\vec{\jmath}\)}
\providecommand{\Oijk}{\(\mathrm{O};\vec{\imath},\vec{\jmath},\vec{k}\)}
\providecommand{\Ouv}{\(\mathrm{O};\vec{u},\vec{v}\)}
\providecommand{\intervFF}[2]{\ensuremath{[#1;#2]}}
\renewcommand{\euro}{\ifmmode\text{\texteuro}\else\texteuro\fi}
\DeclareUnicodeCharacter{00B0}{\ifmmode^{\circ}\else\textdegree\fi}
\DeclareUnicodeCharacter{20AC}{\ifmmode\text{\texteuro}\else\texteuro\fi}
% Énoncés : Éducation nationale. Transcription et corrigés : APMEP, Denis Vergès et contributeurs. Découpage et classement : CoopMaths, Sébastien Lozano, Éric Elter et contributeurs. Adaptation de la présentation : Maths Explorer.
% CC BY-SA : https://coopmaths.fr/www/about/
% https://coopmaths.fr/alea/static/bac/2023/tex/bac_2023_03_sujet2_etrangers_groupe1_4_cor.tex
\begin{document}
\noindent\textit{APMEP / CoopMaths. CC BY-SA. Présentation adaptée par Maths Explorer.}\par\medskip
\bigskip

Vérifions d'abord la condition 1 portant sur la première phase.

Nous avons :
\begin{itemize}
\item une expérience à deux issues (le candidat est qualifié), le succès est \og être qualifié\fg{}, a une probabilité $p = 0,6$ ;
\item cette expérience est répétée quatre fois (quatre candidats) de façon identique et indépendante, donc $n=4$ répétitions;
\item dans ce schéma de Bernoulli de paramètres $n=4$ et $p=0,6$, on note $X$ la variable aléatoire qui compte le nombre de succès.
\end{itemize}

Avec ces éléments, on peut dire que $X$ suit la loi binomiale $ \mathcal{B}\left(4 ~;~ 0,6\right)$.

La condition 1 \og la deuxième phase doit avoir lieu dans 80\% des cas \fg{} correspond à l'évènement \og au moins deux candidats sont sélectionnés \fg{} a une probabilité supérieure ou égale à 0,8 (c'est à dire correspond à $P(X \geqslant 2)\geqslant 0,8$.

Avec la calculatrice : $\displaystyle P(X \geqslant  2)  \approx 0,821$.

On a donc bien : $P(X \geqslant 2)\geqslant 0,8$, la condition 1 est donc vérifiée.

\bigskip 

Pour la condition 2 :

Soit $T$ la variable aléatoire donnant la durée de la deuxième phase.

Les valeurs prises par $T$ sont : 0, 5, 9 et 11.

La deuxième phase dure 11 minutes lorsque quatre candidats sont sélectionnés donc

 $P(T=11)=P(X=4)= \binom{4}{4}\times 0,6^4\times 0,4^{0}=0,6^4= \np{0,1296}$.

De même :

$P(T=9)=P(X=3)= \binom{4}{3}\times 0,6^3\times 0,4^1= \np{0,3456}$.

$P(T=5)=P(X=2)= \binom{4}{2}\times 0,6^2\times 0,4^2= \np{0,3456}$.

$P(T=0)=P(X=0)+P(X=1)= \binom{4}{0}\times 0,6^0\times 0,4^4+ \binom{4}{1}\times 0,6^1\times 0,4^3= \np{0,0256} + \np{0,1536} = \np{0,1792}$


On peut résumer les résultats dans le tableau de la loi de probabilité de la loi $T$ : 

\begin{center}
\begin{tabular}{|l|c|c|c|c|}\hline 
$t_i$ 		& 0 			& 5 			& 9 			& 11\\\hline
$P(T=t_i)$ 	& \np{0,1792} 	& \np{0,3456} 	& \np{0,3456} 	& \np{0,1296}\\\hline
\end{tabular}
\end{center}

La durée moyenne du jeu correspond à l'espérance de la variable aléatoire $T$.

$E(T)= 0 \times P(T=0) + 5 \times P(T=5) + 9 \times P(T=9) + 11 \times P(T=11)$

$E(T) = 0 + 5\times \np{0,3456}+ 9 \times \np{0,3456} + 11 \times \np{0,1296}= 6,264$

La durée moyenne du jeu est donc supérieur à  6 minutes, la deuxième condition n'est pas vérifiée, le jeu ne peut donc pas être retenu.

\end{document}
